\section{Conclusion}
\label{sec:conclusion}

In this paper, we identify the reasoning limitations of ORM-based GRPO on mathematical reasoning tasks. While a simple binary ORM reward is effective, its lack of supervision over reasoning quality causes all reasoning tokens to share the same credit. On the other hand, a naive integration of PRM fails to avoid reward hacking.

To address this challenge, we propose \method{}, which incorporates a rubric-based PRM into GRPO through decoupled advantage normalization. By composing the advantage from independently normalized outcome and process components, with the process advantage normalized only over correct responses, \method{} prevents reward hacking while providing supervision over reasoning quality.

Experiments across four model configurations (3B--14B) and six benchmarks show that \method{} consistently outperforms both GRPO and DAPO baselines, with gains growing as model scale increases. In future work, we aim to explore adaptive weighting between outcome and process signals and extension to broader domains such as scientific reasoning.
